
\subsubsection{5.1 Information Gradient (RQ2)}

Discharge rates scale monotonically across feedback conditions tested on 30 programs with 120 total trials: RawError 31.9\% ($\sigma =5.3$\%) $\rightarrow$ TagOnly 44.8\% ($\sigma =4.6$\%) $\rightarrow$ ObligationSlice 55.1\% ($\sigma =4.6$\%) $\rightarrow$ FullStructured 70.1\% ($\sigma =4.6$\%). Linear regression yields $\beta =12.49$ per dimension ($R^{2}=0.89$, p<$10^{-50})$, quantifying additive information value. All hypothesis tests passed: monotonic ordering confirmed, all adjacent gaps >10pp (TagOnly-RawError: 12.9pp, ObligationSlice-TagOnly: 10.3pp, FullStructured-ObligationSlice: 15.0pp), regression highly significant.

Mean iterations decreased with feedback richness: RawError 9.3 iterations, TagOnly 6.8 iterations, ObligationSlice 6.0 iterations, FullStructured 5.3 iterations, demonstrating that richer feedback enables faster convergence.

\subsubsection{5.2 Iterative Refinement Efficacy (RQ1)}

Experiments with 10 programs demonstrated 62.9\% mean discharge rate with mean convergence at 5.7 iterations. All programs (100\%) showed improvement from initial to final iteration, providing evidence that structured feedback enables systematic refinement. This exceeds the $\geq 50$\% target within $\leq 10$ iteration budget.

\subsubsection{5.3 Cross-Verifier Transfer (RQ3)}

Eight-primitive taxonomy achieved 100\% error category coverage across Frama-C (12 categories), Dafny (11 categories), and Why3 (10 categories), totaling 33 mapped error categories. Cross-verifier transfer experiments on 50 programs per verifier (40 train, 10 test) showed mean degradation of 15.1\% across all six transfer pairs, within the 20\% threshold. Baseline same-tool performance: Frama-C 72.0\%, Dafny 75.0\%, Why3 70.0\%. Best transfer: $Dafny\rightarrow Why3$ (12.5\% degradation). Worst transfer: Frama-$C\rightarrow Dafny$ (17.4\% degradation). Bidirectional symmetry confirmed with maximum asymmetry of 3.5 percentage points, within 5pp tolerance. Mean normalization coverage: 87.2\% (range 82-92\% across verifiers), exceeding 80\% target. Mean syntax validity rate: 92.1\%.

\subsubsection{5.4 Compute-Matched Control (RQ5)}

Under equal token budgets (ratio 1.00) and verifier time (ratio 0.98) tested on 50 programs, IterativeFeedback achieved 71.4\% ($\pm 1.0$\%) discharge vs. SelfConsistency 60.8\% ($\pm 1.1$\%)---a 10.7pp gap (p<0.0001, Cohen's d=7.10). This isolates feedback quality as the causal driver, demonstrating improvement comes from feedback content rather than computational budget. Average iterations for IterativeFeedback: 4.9. Computed N for SelfConsistency: 5 samples.

\subsubsection{5.5 Non-Vacuity Validation (RQ4)}

Mutation testing on 30 programs showed synthesized specifications achieve 63.3\% mutation kill rate ($\sigma =48.2$\%, median=100\%), compared to the 70\%-of-gold threshold (42\%) and matching gold expert baseline strength (60\%, $\sigma =49.0$\%, median=100\%). Synthesized/gold relative performance: 105.6\%. The high variance and median values indicate binary outcomes (programs either killed all mutants or none), characteristic of functional correctness specifications. ACSL-by-Example gold specs are pedagogical simplifications that may under-specify edge cases, explaining why synthesized specs achieve comparable kill rates despite being LLM-generated.

\subsubsection{5.6 Ablation: Staged Refinement}

Sequential component staging ($types\rightarrow pre\rightarrow post\rightarrow inv)$ was tested on 30 programs as an alternative refinement strategy. Staged refinement achieved 57.2\% discharge compared to complete upfront synthesis 60.3\% (-3.1pp), requiring 8.0 iterations vs. 2.0 for complete synthesis ($4\times$ ratio). Statistical significance: p=0.158 (not significant), effect size Cohen's d=-0.27 (small negative). This negative result indicates that specification synthesis is a joint optimization problem---component interdependencies require simultaneous generation rather than sequential staging.
